\documentclass[11pt]{amsart}

\usepackage[T1]{fontenc}
\usepackage[utf8]{inputenc}
\usepackage{amsmath,amssymb,amsthm,mathtools}
\usepackage{stmaryrd}     % \llbracket \rrbracket
\usepackage{xcolor}
\usepackage{etoolbox}
\usepackage{hyperref}
\hypersetup{
  colorlinks = true,
  linkcolor  = blue!60!black,
  citecolor  = blue!60!black,
  urlcolor   = blue!60!black
}
\usepackage{cleveref}

% ---------------------------------------------------------------------
% Macros. Kept in this file rather than in a macros.tex of their own, so
% that the paper is a single file while it is being drafted. Split it off
% when the paper grows sections.
% ---------------------------------------------------------------------

% We quote raw Agda identifiers inline, and these use the same Unicode
% glyphs as the Agda source. Rather than switching to a Unicode engine,
% we teach pdflatex each glyph we use. Extend this list whenever a fresh
% identifier introduces a new character; compiling fails loudly with
% "Unicode character ... not set up" if one is missing.
\usepackage{newunicodechar}
\newunicodechar{¬}{\ensuremath{\neg}}
\newunicodechar{×}{\ensuremath{\times}}
\newunicodechar{Γ}{\ensuremath{\Gamma}}
\newunicodechar{Δ}{\ensuremath{\Delta}}
\newunicodechar{Π}{\ensuremath{\Pi}}
\newunicodechar{Σ}{\ensuremath{\Sigma}}
\newunicodechar{Ω}{\ensuremath{\Omega}}
\newunicodechar{α}{\ensuremath{\alpha}}
\newunicodechar{β}{\ensuremath{\beta}}
\newunicodechar{δ}{\ensuremath{\delta}}
\newunicodechar{ε}{\ensuremath{\varepsilon}}
\newunicodechar{η}{\ensuremath{\eta}}
\newunicodechar{θ}{\ensuremath{\theta}}
\newunicodechar{ι}{\ensuremath{\iota}}
\newunicodechar{μ}{\ensuremath{\mu}}
\newunicodechar{ν}{\ensuremath{\nu}}
\newunicodechar{ρ}{\ensuremath{\rho}}
\newunicodechar{σ}{\ensuremath{\sigma}}
\newunicodechar{τ}{\ensuremath{\tau}}
\newunicodechar{φ}{\ensuremath{\varphi}}
\newunicodechar{χ}{\ensuremath{\chi}}
\newunicodechar{ω}{\ensuremath{\omega}}
\newunicodechar{ᵀ}{\ensuremath{^{\mathrm{T}}}}
\newunicodechar{ᶜ}{\ensuremath{^{\mathrm{c}}}}
\newunicodechar{′}{\ensuremath{'}}
\newunicodechar{⁻}{\ensuremath{^{-}}}
\newunicodechar{⁺}{\ensuremath{^{+}}}
\newunicodechar{¹}{\ensuremath{^{1}}}
\newunicodechar{₀}{\ensuremath{_{0}}}
\newunicodechar{₁}{\ensuremath{_{1}}}
\newunicodechar{₂}{\ensuremath{_{2}}}
\newunicodechar{₃}{\ensuremath{_{3}}}
\newunicodechar{𝓑}{\ensuremath{\mathcal{B}}}
\newunicodechar{𝓞}{\ensuremath{\mathcal{O}}}
\newunicodechar{𝓤}{\ensuremath{\mathcal{U}}}
\newunicodechar{ℓ}{\ensuremath{\ell}}
\newunicodechar{ℕ}{\ensuremath{\mathbb{N}}}
\newunicodechar{𝟘}{\ensuremath{\mathbf{0}}}
\newunicodechar{𝟙}{\ensuremath{\mathbf{1}}}
\newunicodechar{𝟚}{\ensuremath{\mathbf{2}}}
\newunicodechar{⋆}{\ensuremath{\star}}
\newunicodechar{→}{\ensuremath{\rightarrow}}
\newunicodechar{⇒}{\ensuremath{\Rightarrow}}
\newunicodechar{↦}{\ensuremath{\mapsto}}
\newunicodechar{∘}{\ensuremath{\circ}}
\newunicodechar{∼}{\ensuremath{\sim}}
\newunicodechar{≃}{\ensuremath{\simeq}}
\newunicodechar{≤}{\ensuremath{\leq}}
\newunicodechar{≥}{\ensuremath{\geq}}
\newunicodechar{＝}{\ensuremath{=}}
\newunicodechar{⊕}{\ensuremath{\oplus}}
\newunicodechar{⊗}{\ensuremath{\otimes}}
\newunicodechar{⊑}{\ensuremath{\sqsubseteq}}
\newunicodechar{⊴}{\ensuremath{\trianglelefteq}}
\newunicodechar{⊲}{\ensuremath{\vartriangleleft}}
\newunicodechar{⋖}{\ensuremath{\lessdot}}
\newunicodechar{⊙}{\ensuremath{\odot}}
\newunicodechar{▸}{\ensuremath{\blacktriangleright}}
\newunicodechar{∋}{\ensuremath{\ni}}
\newunicodechar{⟦}{\ensuremath{\llbracket}}
\newunicodechar{⟧}{\ensuremath{\rrbracket}}
\newunicodechar{⟨}{\ensuremath{\langle}}
\newunicodechar{⟩}{\ensuremath{\rangle}}

% ---------------------------------------------------------------------
% Basic types and system T
\newcommand{\N}{\mathbb{N}}
\newcommand{\Baire}{\N^{\N}}
% Dialogue trees are written B, as in the formalization, and Brouwer
% codes are written with a calligraphic B, which is the renaming the
% Agda files make for the same reason.
\newcommand{\Dial}{\mathrm{B}}
\newcommand{\Tm}{\mathrm{T}}
\newcommand{\TmO}{\mathrm{T}^{\Omega}}
\newcommand{\gr}{\iota}                      % the ground type
\newcommand{\arr}{\Rightarrow}
\newcommand{\sem}[1]{\llbracket #1 \rrbracket}
\newcommand{\dsem}[1]{\llbracket #1 \rrbracket_{\Dial}}
\newcommand{\dtree}{\mathrm{dialogue\text{-}tree}}
\newcommand{\height}{\mathrm{height}}
\newcommand{\kext}[1]{#1^{*}}                % Kleisli extension
\newcommand{\iterop}{\mathrm{iter}}
% The combinators are set in bold upright, the constructors of the
% Brouwer codes in sans serif, so that the combinator S and the
% successor of a code are told apart on sight.
\newcommand{\Iterc}{\mathbf{Iter}}
\newcommand{\Kc}{\mathbf{K}}
\newcommand{\Sc}{\mathbf{S}}
\newcommand{\Zeroc}{\mathbf{Zero}}
\newcommand{\Succc}{\mathbf{Succ}}

% ---------------------------------------------------------------------
% Brouwer codes and their arithmetic. The operations on codes are
% written with the ordinary symbols of ordinal arithmetic, which is what
% they denote, and the macros survive so that the source says which is
% meant.
\newcommand{\Brouwer}{\mathcal{B}}
\newcommand{\bZ}{\mathsf{Z}}
\newcommand{\bS}{\mathsf{S}}
\newcommand{\bL}{\mathsf{L}}
\newcommand{\bplus}{\mathbin{+}}
\newcommand{\btimes}{\mathbin{\cdot}}
\newcommand{\num}[1]{\underline{#1}}          % the code of a numeral
\newcommand{\eps}{\varepsilon_{0}}
\newcommand{\tower}[1]{\omega_{#1}}           % the tower of exponentials

% ---------------------------------------------------------------------
% Majorants and data
\newcommand{\Maj}{\mathrm{Maj}}
\newcommand{\maj}{\mathrel{\lhd}}             % the majorization relation R
\newcommand{\Mj}{\mu}                         % the majorant of a term
\newcommand{\ord}{\mathrm{ord}}
% \wd is a TeX primitive, so the width of the data of a type is \dwd.
\newcommand{\dwd}{\mathrm{w}}
\newcommand{\W}{\mathrm{W}}
\newcommand{\tsum}[1]{\textstyle\sum #1}
\newcommand{\tle}{\mathrel{\sqsubseteq}}      % componentwise order on data
\newcommand{\tsub}{\mathrel{\trianglelefteq}} % a datum inside a wider one
\newcommand{\dom}{\mathrel{\preccurlyeq}}     % the domination predicate
\newcommand{\act}[2]{#1[#2]}                  % action of a cascade or a map
\newcommand{\rd}[1]{\langle #1 \rangle}       % the map a tuple is read as

% ---------------------------------------------------------------------
% Pointers into the Agda formalization are set in a light blue so that
% they are easy to tell apart from the mathematical text. They can be
% turned off in one place by setting the switch to \agdarefsfalse, which
% makes every \AgdaRefs print nothing while leaving its argument in the
% source. They are on while the paper is being verified.
\definecolor{agdablue}{HTML}{5E8CB8}
\newif\ifagdarefs
\agdarefstrue
\newcommand{\AgdaRefs}[1]{\ifagdarefs{\footnotesize\color{agdablue}(#1)}\else\unskip\fi}

% Claims made in running prose, outside any numbered environment, are
% named with \claim{name} so that a companion file can point back at
% them, and the names are written to \jobname.claims in the order they
% occur. Every claim of this paper sits in a numbered environment with a
% pointer of its own, so nothing uses this yet.
\newwrite\claimfile
\immediate\openout\claimfile=\jobname.claims
\newcommand{\claim}[1]{%
  \ifagdarefs{\footnotesize\color{agdablue}[#1]}\else\unskip\fi
  \immediate\write\claimfile{#1}}

%   \Agda{Module.Path}{identifier-or-identifiers}
% A module-only pointer is written with an empty second argument, and the
% separating dot is then suppressed.
\newcommand{\Agda}[2]{{\footnotesize\color{agdablue}\textsf{#1}%
  \ifblank{#2}{}{\texttt{.}\allowbreak\texttt{#2}}}}
\newcommand{\AgdaModule}[1]{{\footnotesize\color{agdablue}\textsf{#1}}}

% Continuation entries in a pointer list are written as bare \texttt{},
% so colour and shrink that too. Every \texttt in this paper is an Agda
% identifier or an Agda compiler flag.
\makeatletter
\let\ttt@plain\texttt
\renewcommand{\texttt}[1]{{\footnotesize\color{agdablue}\ttt@plain{#1}}}
\makeatother

% ---------------------------------------------------------------------
% All theorem-like environments share the theorem counter, so that the
% numbering is a single sequence per section. Sharing a counter the plain
% way would make cleveref call every one of them a theorem, since it
% names references by counter, so aliascnt gives each environment its own
% counter name kept in step with the theorem counter.
\usepackage{aliascnt}
\newcommand{\sharedtheorem}[3]{%
  \newaliascnt{#1}{theorem}%
  \newtheorem{#1}[#1]{#2}%
  \aliascntresetthe{#1}%
  \crefname{#1}{#2}{#3}%
  \Crefname{#1}{#2}{#3}%
}

\theoremstyle{plain}
\newtheorem{theorem}{Theorem}[section]
\crefname{theorem}{Theorem}{Theorems}
\Crefname{theorem}{Theorem}{Theorems}
\crefname{section}{Section}{Sections}
\Crefname{section}{Section}{Sections}
\crefname{appendix}{Appendix}{Appendices}
\Crefname{appendix}{Appendix}{Appendices}
\crefname{subsection}{Section}{Sections}
\Crefname{subsection}{Section}{Sections}
\sharedtheorem{proposition}{Proposition}{Propositions}
\sharedtheorem{lemma}{Lemma}{Lemmas}
\sharedtheorem{corollary}{Corollary}{Corollaries}

\theoremstyle{definition}
\sharedtheorem{definition}{Definition}{Definitions}
\sharedtheorem{example}{Example}{Examples}
\sharedtheorem{examples}{Examples}{Examples}
\sharedtheorem{remark}{Remark}{Remarks}
\sharedtheorem{notation}{Notation}{Notations}
\sharedtheorem{question}{Question}{Questions}

% ---------------------------------------------------------------------
% The short form in brackets is the running head, the full title being
% too long for one line at the top of a page.
\title[Dialogue trees of system T]{The dialogue trees of system T have
  height below $\varepsilon_0$, and this bound is sharp}
\date{\today}

\begin{document}

\begin{abstract}
  The dialogue interpretation of G\"odel's system T assigns to each
  closed term of type $(\gr\arr\gr)\arr\gr$ a well founded tree,
  branching over the natural numbers, whose nodes record the queries
  the term makes to its oracle. We prove that the height of this tree,
  measured as a Brouwer ordinal code, is below~$\eps$, and that
  $\eps$ is the least such bound, a closed term attaining the height
  $\tower n$ for every~$n$. The upper bound assigns to every term,
  by induction on its type, a majorant built from addition, suprema
  and suprema of iterates, and bounds the majorant by a tuple of codes
  graded by the order of the type, in which an entry is raised only by
  entries of its own order or higher. The lower bound runs Church
  encodings of the approximants of~$\eps$ against the oracle. The
  development is formalized in Agda, with function extensionality as
  its only assumption. Xiao and Sun have proved the same bound, with
  a computable tower level for each term, over classical ordinals.
  The proof offered here is an alternative one, constructive, and it
  settles the sharpness question that they leave open.
\end{abstract}

\maketitle

\noindent
This paper and its accompanying formalization were developed by LLM's
under Mart\'in Escard\'o's direction. He has verified the correctness
of the formulation of the main theorems in the Agda code and takes
responsibility for this. However, he takes no responsibility for the
correctness of this paper, but see the precautions taken
in~\cref{sec:appendix}. This paper would greatly benefit from further
human care.

\section{Introduction}
\label{sec:intro}

G\"odel's system T defines only continuous functionals
$\Baire\to\N$. One proof of this, by effectful forcing, interprets
each closed term $t$ of type $(\gr\arr\gr)\arr\gr$ as a dialogue
tree~\cite{Escardo2013}: a well founded tree, branching over the
natural numbers, whose nodes carry the queries the term makes to its
oracle and whose leaves carry its answers. Continuity is read off the
tree, since a leaf is reached after finitely many queries. The trees
are Church encoded within system T itself
in~\cite{EscardoPaivaRahliTosun2025}.

The tree is well founded, so it has an ordinal height, and the
question is how large that height can be. The answer is that the
height is below~$\eps$, and that no smaller ordinal bounds the heights
of all the trees.

The first of these was proved by Xiao and Sun~\cite{XiaoSun2026}, who
compute from each term a tower level below~$\eps$ bounding the height
of its tree, by translating recursors into closed infinitary templates
and removing redexes in a sequence of passes indexed by type level.
Their ordinals are the classical ones, and their formalization assumes
excluded middle for the propositions of the second universe, along
with univalence, propositional truncation and set replacement. They
write that they do not claim the bound sharp, and they leave open
whether the heights are cofinal in~$\eps$.

What this paper adds is twofold. The proof is constructive: the
ordinals are Brouwer codes with their syntactic order, so the whole
development assumes function extensionality and nothing else, and a
bound in that order implies the classical one. And the heights are
shown to be cofinal in~$\eps$ (\cref{cor:least}), which answers the
question Xiao and Sun leave open and makes $\eps$ the least bound.
What this paper does not supply is their computable tower level: the
bound obtained here is $\eps$ itself, with no finite level named for a
given term.

Two features of the problem shape the proof. The first is that the
height of the tree of a term is not a function of the values the term
computes. A term whose numerical output is small may ask the oracle a
number of questions that no bound on its output predicts, so an
ordinal assignment bounding the values of a term, such as Howard's
\cite{Howard1970}, does not by itself bound the height of its dialogue
tree. The second is that the recursor at higher types iterates a
functional rather than a number, and the bound for the iterates has to
be hereditary in the type.

We meet the first by bounding the height directly, with a logical
relation over the types of system T whose ground case reads
\[
  \height(d) \le a ,
\]
for a dialogue tree $d$ and a Brouwer code~$a$. We meet the second by
making the majorant hereditary, with a pointwise supremum and a
pointwise translation at every type, so that the recursor acquires a
majorant at every type and not only at the ground type. This reduces
the theorem to a statement of ordinal arithmetic
(\cref{thm:reduction}): the majorant of a closed term of ground type
is below~$\eps$.

A majorant is a functional over Brouwer codes, built from addition
and from suprema, among them the suprema of iterates. To bound one by
an ordinal below~$\eps$ we bound it by a tuple of codes graded by the
order of the type: a
translation at order~$0$, an exponent at order~$1$, the exponent of
that exponent at order~$2$, and so on. One rule governs the whole
construction, that an entry is raised only by entries of its own order
or higher, which is the rule that iterating $x\mapsto\omega^{x}$
breaks. The maps between such tuples that obey the rule form a
calculus with application, abstraction and projection
(\cref{sec:graded}), and the two hard clauses of the induction on
terms, the one for~$\Sc$ and the one for the recursor, are then
written as a reader of $\lambda f\,g\,a.\,f\,a\,(g\,a)$ would write
them.

For the lower bound we write Church encodings of the approximants
of~$\eps$ as terms of system T, after Coquand, Hancock and
Setzer~\cite{CoquandHancockSetzer1997}, and run them against the
oracle through an algebra that places one query above the tree built
so far. A parametricity argument shows that the resulting tree has the
intended height.

\subsection*{The formalization}

Every result below carries a pointer, set in blue, to the Agda
definition that proves it. The modules named live under the prefix
\AgdaModule{DialogueTreesHeight}, which the pointers leave out,
within the development accompanying this paper, available at
\url{https://github.com/martinescardo/DialogueTreesHeight}, and the
modules whose names begin with \AgdaModule{Ordinals} or \AgdaModule{EffectfulForcing} are those
of the \texttt{TypeTopology} library. Everything is checked with the
flags \texttt{--safe} and \texttt{--without-K}, with no postulates and
no holes. The one assumption is function extensionality, which is a
parameter of every module and is used at the limit case of each
arithmetic law proved by induction. The pointers can be suppressed by
setting \texttt{\textbackslash agdarefsfalse} in the preamble.

\section{Dialogue trees and their heights}
\label{sec:dialogue}

We recall the dialogue trees of~\cite{Escardo2013}, the functor and
monad they form, and the dialogue interpretation that assigns a tree
to each closed term of system T.

\begin{definition}[Dialogue trees]
\label{def:dialogue-tree-type}
For a type~$X$, the type $\Dial X$ of dialogue trees over~$X$ is
inductively generated by
\[
  \eta : X \to \Dial X , \qquad
  \beta : (\N \to \Dial X) \to \N \to \Dial X .
\]
A tree $\beta\,\varphi\,i$ asks the oracle for its value at~$i$ and
continues with $\varphi\,j$ when the answer is~$j$, and a tree $\eta\,x$
answers~$x$. Running a tree against an oracle $\alpha : \Baire$ gives
\[
  \mathrm{dialogue} : \Dial X \to \Baire \to X ,
\]
defined by $\mathrm{dialogue}\,(\eta\,x)\,\alpha = x$ and
$\mathrm{dialogue}\,(\beta\,\varphi\,i)\,\alpha =
 \mathrm{dialogue}\,(\varphi\,(\alpha\,i))\,\alpha$.
\AgdaRefs{\Agda{EffectfulForcing.MFPSAndVariations.Dialogue}{D};
\texttt{B}, \texttt{dialogue}.}
\end{definition}

The operation $\Dial$ is a functor, the tree $\Dial f\,d$ being $d$
with its leaves relabelled by~$f$, and a monad. We write $\kext g$ for
the Kleisli extension of $g : X \to \Dial Y$, which grafts the tree
$g\,x$ onto each leaf $\eta\,x$, and
$\mathrm{generic} = \kext{(\beta\,\eta)}$ for the generic oracle,
which asks the oracle at the value a tree computes.
\AgdaRefs{\Agda{EffectfulForcing.MFPSAndVariations.Dialogue}{kleisli-extension};
\texttt{generic}, \texttt{B-functor}.}

We write $\iterop$ for iteration on the natural numbers,
$\iterop\,f\,x\,0 = x$ and
$\iterop\,f\,x\,(n+1) = f\,(\iterop\,f\,x\,n)$.
\AgdaRefs{\Agda{EffectfulForcing.MFPSAndVariations.Combinators}{iter}.}

\begin{definition}[The dialogue interpretation]
\label{def:dialogue-interpretation}
Write $\Tm\sigma$ for the closed terms of the combinatory system T,
with constants $\Zeroc$, $\Succc$, $\Iterc$, $\Kc$ and~$\Sc$, and
$\TmO\sigma$ for the same language extended with a constant $\Omega$
of type $\gr\arr\gr$ for the oracle. The dialogue interpretation of a
type is
\[
  \dsem{\gr} = \Dial\N , \qquad
  \dsem{\sigma\arr\tau} = \dsem{\sigma} \to \dsem{\tau} ,
\]
and the dialogue interpretation of a closed term $t$ of $\TmO\sigma$
is an element $\dsem{t}$ of~$\dsem\sigma$, the same brackets serving
for both because the argument says which is meant. It interprets
$\Omega$ as $\mathrm{generic}$, $\Zeroc$ as $\eta\,0$, $\Succc$ as
$\Dial$ applied to the successor, $\Kc$ and~$\Sc$ as the projection
and the substitution, and the recursor $\Iterc$ at type~$\sigma$ as
\[
  \iterop'\,f\,x = \kext{(\iterop\,f\,x)} ,
\]
the Kleisli extension taken at type~$\sigma$. The dialogue tree of a
closed term $t : (\gr\arr\gr)\arr\gr$ is
$\dtree(t) = \dsem{\overline t \cdot \Omega}$, where $\overline t$ is
the image of~$t$ in~$\TmO$.
\AgdaRefs{\Agda{EffectfulForcing.MFPSAndVariations.MFPS-XXIX}{B⟦\_⟧};
\texttt{B-Set⟦\_⟧}, \texttt{Kleisli-extension}, \texttt{iter'},
\texttt{dialogue-tree};
\Agda{EffectfulForcing.MFPSAndVariations.CombinatoryT}{T},
\texttt{TΩ}, \texttt{embedding}.}
\end{definition}

The Kleisli extension in that definition is taken at an arbitrary
type, which asks for more than the monad supplies. For
$g : \N \to \dsem\sigma$ and a tree $d$, define $\kext g\,d$ in
$\dsem\sigma$ by recursion on the type: at the ground type it is the
Kleisli extension of the monad, and at an arrow it acts under the
argument,
\[
  (\kext g\,d)\,s = \kext{(k \mapsto g\,k\,s)}\,d .
\]
\AgdaRefs{\Agda{EffectfulForcing.MFPSAndVariations.MFPS-XXIX}{Kleisli-extension}.}

The constant $\Iterc$ iterates a function rather than performing
primitive recursion, and we call it the recursor throughout. We also
abbreviate combinatory terms by $\lambda$-notation, as in
$\lambda f.\,f\,0$ for $\Sc\cdot I\cdot(\Kc\cdot\Zeroc)$, where
$I = \Sc\cdot\Kc\cdot\Kc$ is the identity, the abbreviation standing
for the term that bracket abstraction produces.

The height of a dialogue tree is defined in the next section, once the
Brouwer codes are at hand.

\section{Brouwer codes and their arithmetic}
\label{sec:codes}

We take the ordinals to be Brouwer codes with their syntactic order,
for which the successor is monotone by construction, define addition,
multiplication, exponentiation and the tower whose supremum is~$\eps$,
record the closure properties of~$\eps$ that the proof uses, and
define the height of a dialogue tree with its three properties.

A classical treatment would interpret the height of a dialogue tree as
an ordinal in the sense of the univalent library, and then the grafting
and generic lemmas need excluded middle, because the successor map on
ordinals is not monotone without it. Brouwer codes with their
syntactic order avoid this.
\AgdaRefs{\Agda{Ordinals.AdditionProperties}{succ-not-necessarily-monotone}.}

\begin{definition}[Brouwer codes and their order]
\label{def:codes}
The type $\Brouwer$ of Brouwer ordinal codes is inductively generated
by
\[
  \bZ : \Brouwer , \qquad
  \bS : \Brouwer \to \Brouwer , \qquad
  \bL : (\N \to \Brouwer) \to \Brouwer .
\]
The order $\le$ on~$\Brouwer$ is inductively generated by four rules,
the least element property, the successor rule, the cocone and the
limit rule, saying that $\bZ$ is least, that the successor is
monotone, that $\bL f$ is above each $f\,n$, and that $\bL f$ is below
every upper bound of the family~$f$:
\[
  \frac{\ }{\bZ \le b} \qquad
  \frac{a \le b}{\bS a \le \bS b} \qquad
  \frac{a \le f\,n}{a \le \bL f} \qquad
  \frac{f\,n \le b \ \text{for every $n$}}{\bL f \le b} .
\]
We write $a < b$ for $\bS a \le b$.
\AgdaRefs{\Agda{Ordinals.BrouwerCodes}{B};
\Agda{BrouwerOrdinals.Order}{\_≤\_},
\Agda{BrouwerOrdinals.Epsilon0}{\_<\_}.}
\end{definition}

A limit is therefore the least upper bound of its family, and the
successor is monotone for arbitrary codes, which is what the classical
reading lacks. The order is reflexive and transitive, transitivity by a
lexicographic recursion that descends on the left proof when it is a
limit and on the right proof otherwise.
\AgdaRefs{\Agda{BrouwerOrdinals.Order}{≤-refl}; \texttt{≤-trans},
\texttt{≤-L-mono}, \texttt{≤-L-upper-bound}.}

\begin{remark}
\label{rem:one-branch}
Two asymmetries of the order govern what can be proved about it. A
lower bound $a \le \bL f$ holds as soon as $a \le f\,n$ for a single
branch~$n$, while an upper bound $\bL f \le b$ asks for every member of
the family. A numeral therefore fits under one branch of a tree, which
is what makes a lower bound of~$\omega$ available by induction
(\cref{prop:omega-attained}), and a limit does not, which is why the
lower bounds of~\cref{sec:sharp} go through terms rather than through
the order alone.
\end{remark}

\begin{definition}[Arithmetic of codes]
\label{def:arithmetic}
Addition, multiplication and base~$\omega$ exponentiation are defined
by recursion on the right argument, respectively on the exponent:
\begin{align*}
  a \bplus \bZ    &= a ,                      &
  a \btimes \bZ   &= \bZ ,                    &
  \omega^{\bZ}    &= \bS\,\bZ ,               \\
  a \bplus \bS b  &= \bS (a \bplus b) ,       &
  a \btimes \bS b &= (a \btimes b) \bplus a , &
  \omega^{\bS a}  &= \omega^{a} \btimes \omega , \\
  a \bplus \bL f  &= \bL (n \mapsto a \bplus f\,n) , &
  a \btimes \bL f &= \bL (n \mapsto a \btimes f\,n) , &
  \omega^{\bL f}  &= \bL (n \mapsto \omega^{f\,n}) .
\end{align*}
The numerals are $\num 0 = \bZ$ and $\num{n+1} = \bS\,\num n$, the
limit $\omega$ is $\bL(n\mapsto\num n)$, the tower is
$\tower 0 = \omega$ and $\tower{n+1} = \omega^{\tower n}$, and
\[
  \eps = \bL(n \mapsto \tower n) .
\]
\AgdaRefs{\Agda{BrouwerOrdinals.Order}{\_⊕\_};
\Agda{BrouwerOrdinals.Multiplication}{ι[\_]}, \texttt{ω}, \texttt{\_⊗\_};
\Agda{BrouwerOrdinals.Epsilon0}{ω\^{}\_}, \texttt{tower}, \texttt{ε₀}.}
\end{definition}

The equation $a \bplus \bS b = \bS(a \bplus b)$ makes addition on the
left commute with successors and limits definitionally, which is what
the inductions below rely on. Addition is associative and monotone in
each argument, multiplication by a fixed code on the left distributes
over addition and is associative and monotone in each argument, and
exponentiation is monotone and inflationary.
\AgdaRefs{\Agda{BrouwerOrdinals.Multiplication}{⊕-assoc};
\texttt{⊕-mono-left}, \texttt{⊗-left-distrib}, \texttt{⊗-assoc};
\Agda{BrouwerOrdinals.Epsilon0}{⊕-mono-right}, \texttt{⊗-mono-left},
\texttt{⊗-mono-right}, \texttt{ω\^{}-mono}, \texttt{≤-ω\^{}}.}

\begin{lemma}[The tower is strictly increasing]
\label{lem:tower}
For every~$n$ we have $\tower n < \tower{n+1}$, and hence
$\tower n < \eps$.
\AgdaRefs{\Agda{BrouwerOrdinals.Epsilon0}{tower-strict};
\texttt{tower-<-ε₀}.}
\end{lemma}

\begin{proof}
A direct induction on~$n$ has to put the successor of the limit
$\tower n$ below the limit $\tower{n+1}$, which by the cocone rule
needs a single member of the latter above the whole of the former,
and the induction hypothesis supplies none. The obstruction is
removed by a descent:
strictness at $\omega^{b}$ reduces, through
$\bS(\omega^{b}) \le \omega^{b} \bplus \omega^{b} \le \omega^{b}
 \btimes \omega$ and monotonicity of the exponential, to strictness
at~$b$, and the descent stops at $\omega < \omega^{\omega}$, which is
proved from $\omega^{\num 1} = \omega$.
\AgdaRefs{\Agda{BrouwerOrdinals.Epsilon0}{S-ω\^{}-step};
\texttt{ω\^{}-ι1}, \texttt{ω<ω\^{}ω}.}
\end{proof}

\begin{proposition}[Closure of $\eps$]
\label{prop:closure}
If $a < \eps$ then $\bS a < \eps$ and $\omega^{a} < \eps$, and if
$b < \eps$ as well then $a \bplus b < \eps$ and $a \btimes b < \eps$.
\AgdaRefs{\Agda{BrouwerOrdinals.Epsilon0}{⊕-<-ε₀}; \texttt{S-<-ε₀},
\texttt{ω\^{}-<-ε₀}; \Agda{BrouwerOrdinals.MultClosure}{⊗-<-ε₀}.}
\end{proposition}

\begin{proof}
A proof of $a < \eps$ can be inverted to a tower level by a single
pattern match, since the cocone rule is the only one that puts a
successor below the limit~$\eps$. Two codes below~$\eps$ are therefore
pushed to a common tower level, and the sum is absorbed one level up
by $\tower n \bplus \tower n \le \tower{n+1}$, the product by
$\tower n \btimes \tower n \le \tower{n+1}$.
\AgdaRefs{\Agda{BrouwerOrdinals.Epsilon0}{tower-double};
\texttt{tower-mono-ℕ}; \Agda{BrouwerOrdinals.MultClosure}{tower-sq},
\texttt{ω\^{}⊗}.}
\end{proof}

\begin{definition}[The height of a dialogue tree]
\label{def:height}
The height of a dialogue tree is the Brouwer code
\[
  \height(\eta\,n) = \bZ , \qquad
  \height(\beta\,\varphi\,i) = \bL(n \mapsto \bS\,\height(\varphi\,n)) ,
\]
that is, $0$ at a leaf and the supremum of the heights of the subtrees
raised by one at a query node. The query index is ignored, the height
depending on the branching structure alone.
\AgdaRefs{\Agda{DialogueTrees.Constructive}{height}.}
\end{definition}

\begin{lemma}[Relabelling, grafting and the generic oracle]
\label{lem:height-three}
Let $f : \N\to\N$, let $g : \N \to \Dial\N$, let $b : \Brouwer$ and let
$d : \Dial\N$.
\begin{enumerate}
\item\label{item:relabel}
  $\height(\Dial f\,d) = \height(d)$.
\item\label{item:graft}
  If $\height(g\,k) \le b$ for every~$k$, then
  $\height(\kext g\,d) \le b \bplus \height(d)$.
\item\label{item:generic}
  $\height(\mathrm{generic}\,d) \le \bS\,\height(d)$.
\end{enumerate}
\AgdaRefs{\Agda{DialogueTrees.Constructive}{height-relabel};
\texttt{height-kleisli-extension}, \texttt{height-generic}.}
\end{lemma}

\begin{proof}
All three are inductions on~$d$. In
(\labelcref{item:relabel}) the height ignores the values at the
leaves, and this is the only one of the three that uses function
extensionality, to pass from a pointwise identification of two
families to an identification of their limits.
In (\labelcref{item:graft}) the summand~$b$ is on the left because the
structure of~$d$ lies above the grafted trees, and the successor step
is unconditional, which is what makes the proof constructive. Part
(\labelcref{item:generic}) amounts to grafting a single query, since
$\mathrm{generic} = \kext{(\beta\,\eta)}$, and is proved by an
induction of its own rather than as an instance
of~(\labelcref{item:graft}).
\end{proof}

\section{Majorants and the reduction to ordinal arithmetic}
\label{sec:majorant}

We assign to each term a hereditary majorant, extend the assignment to
the recursor at every type by equipping the majorants with a pointwise
supremum and a pointwise translation, and reduce the theorem to the
statement that the majorant of a closed term of ground type is
below~$\eps$.

\begin{definition}[Majorants]
\label{def:majorant}
The type of majorants of a system T type is defined by recursion on
the type,
\[
  \Maj(\gr) = \Brouwer , \qquad
  \Maj(\sigma\arr\tau) = \Maj(\sigma) \to \Maj(\tau) ,
\]
and so is the relation of being majorized:
\begin{align*}
  x \maj_{\gr} a &\quad\text{means}\quad \height(x) \le a , \\
  F \maj_{\sigma\arr\tau} \varphi &\quad\text{means that}\quad
    x \maj_{\sigma} a \ \text{implies}\ F\,x \maj_{\tau} \varphi\,a ,
\end{align*}
the latter for all $x$ and~$a$.
\AgdaRefs{\Agda{DialogueTrees.Majorant}{Maj}; \texttt{R}.}
\end{definition}

\begin{definition}[Majorants of the combinators]
\label{def:majorant-combinators}
The constant $\Zeroc$ has majorant~$\bZ$, the successor has the
identity, since relabelling preserves the height, the oracle has the
successor $a \mapsto \bS a$, and $\Kc$ and~$\Sc$ have the projection
$(a,b)\mapsto a$ and the substitution
$(\varphi,\gamma,a) \mapsto \varphi\,a\,(\gamma\,a)$. Each combinator
is majorized by its majorant.
\AgdaRefs{\Agda{DialogueTrees.Majorant}{μ-Zero}; \texttt{μ-Succ},
\texttt{μ-Ω}, \texttt{μ-K}, \texttt{μ-S}, \texttt{R-Zero},
\texttt{R-Succ}, \texttt{R-Ω}, \texttt{R-K}, \texttt{R-S}.}
\end{definition}

The recursor is the remaining constant. At ground type its majorant is
the supremum of the iterates of a family of codes, and at a type
$\sigma\arr\tau$ a majorant was missing. We supply one at every type.

\begin{definition}[Supremum and translation]
\label{def:maj-sup}
The pointwise supremum and the pointwise right translation on
majorants are defined by recursion on the type, acting under the
argument at an arrow:
\begin{align*}
  \textstyle\sup_{\gr} h &= \bL h , &
  a \bplus_{\gr} \nu &= a \bplus \nu , \\
  \textstyle(\sup_{\sigma\arr\tau} h)\,a
    &= \textstyle\sup_{\tau} (k \mapsto h\,k\,a) , &
  (\varphi \bplus_{\sigma\arr\tau} \nu)\,a
    &= (\varphi\,a) \bplus_{\tau} \nu .
\end{align*}
\AgdaRefs{\Agda{DialogueTrees.MajorantAll}{Maj-sup};
\texttt{Maj-⊕},
\texttt{μ-Kleisli}.}
\end{definition}

\begin{lemma}[The hereditary Kleisli bound]
\label{lem:kleisli-bound}
Let $g : \N \to \dsem\sigma$ and $h : \N \to \Maj(\sigma)$ satisfy
$g\,k \maj_{\sigma} h\,k$ for every~$k$, and let $d$ be a dialogue tree
with $\height(d) \le \nu$. Then
\[
  \kext g\,d \ \maj_{\sigma}\
  \big(\textstyle\sup_{\sigma} h\big) \bplus_{\sigma} \nu .
\]
\AgdaRefs{\Agda{DialogueTrees.MajorantAll}{R-Kleisli}.}
\end{lemma}

\begin{proof}
By induction on~$\sigma$. At the ground type the grafting lemma
applies with the bound $\bL h$, which the upper bound property of the
limit supplies, and monotonicity of addition in its right argument
then passes from $\height(d)$ to~$\nu$; see
\cref{lem:height-three}(\labelcref{item:graft}). The arrow case holds
by definition on both
sides, the Kleisli extension, the supremum and the translation all
passing under the argument in the same way.
\end{proof}

\begin{definition}[The majorant of a term]
\label{def:majorant-term}
The majorant of the recursor at type~$\sigma$ is
\[
  \Mj_{\Iterc}\,\varphi\,a\,\nu
  = \big(\textstyle\sup_{\sigma}(k \mapsto \iterop\,\varphi\,a\,k)\big)
     \bplus_{\sigma} \nu ,
\]
and the majorant $\Mj(t)$ of a term of~$\TmO$ is defined by the
constants above and $\Mj(t\cdot u) = \Mj(t)\,\Mj(u)$.
\AgdaRefs{\Agda{DialogueTrees.MajorantAll}{μ-Iterᵀ}; \texttt{μᵀ},
\texttt{R-Iterᵀ}.}
\end{definition}

\begin{theorem}[The fundamental theorem for majorants]
\label{thm:fundamental-majorant}
Every closed term $t$ of~$\TmO\sigma$ satisfies
$\dsem t \maj_{\sigma} \Mj(t)$. In particular, for a closed term
$t : (\gr\arr\gr)\arr\gr$ of system T,
\[
  \height(\dtree(t)) \le \Mj(\overline t \cdot \Omega) .
\]
\AgdaRefs{\Agda{DialogueTrees.MajorantAll}{fundamentalᵀ};
\texttt{dialogue-tree-height-≤-μᵀ}.}
\end{theorem}

\begin{proof}
By induction on the term. Every case but the recursor is the majorant
of the corresponding combinator, and the recursor is
\cref{lem:kleisli-bound} applied to the iterates, which are majorized
by induction on the iteration count.
\end{proof}

\begin{theorem}[The reduction]
\label{thm:reduction}
Let $t : (\gr\arr\gr)\arr\gr$ be a closed term of system T. If
$\Mj(\overline t \cdot \Omega) < \eps$ then
$\height(\dtree(t)) < \eps$.
\AgdaRefs{\Agda{DialogueTrees.MajorantAll}{dialogue-height-<-ε₀}.}
\end{theorem}

\begin{proof}
\Cref{thm:fundamental-majorant} bounds the height by the majorant, and
the successor is monotone, so
\[
  \bS\,\height(\dtree(t)) \ \le\ \bS\,\Mj(\overline t \cdot \Omega)
  \ \le\ \eps ,
\]
the second inequality being the hypothesis.
\end{proof}

The hypothesis concerns Brouwer codes alone, $\Mj(t)$ being a
hereditary functional over codes built from addition, suprema and
suprema of iterates. The rest of the paper proves that hypothesis, in
the form $\Mj(t) < \eps$ for every closed term $t$ of~$\TmO\gr$.

\section{Graded maps}
\label{sec:graded}

We define the data and their grading, the scalar maps between them
whose exponents are again maps, and the rectangular maps between data
of any two widths, with the four facts about them that the induction
on terms uses.

\subsection{Data and their grading}
\label{sec:data}

A \emph{datum of width~$n$} is an $n$-tuple of Brouwer codes, written
$\Brouwer^{n}$. Its entry at position~$j$ is its entry \emph{of
order~$j$}, where order~$0$ is a translation, order~$1$ an exponent,
order~$2$ the exponent of an exponent, and so on. We write $\tsum t$
for the sum of the entries of~$t$, we write $t \tle u$ for the
componentwise order on data of equal width, and we say that a datum is
\emph{small} when each of its entries is below~$\eps$. A datum $t$ of
width~$k$ sits inside a datum $u$ of width~$n$, written $t \tsub u$,
when $k \le n$ and each entry of~$t$ is below the entry of~$u$ of the
same order.
\AgdaRefs{\Agda{BrouwerOrdinals.MoveTuple}{Tup}; \texttt{tsum},
\texttt{\_⊑\_}, \texttt{All<}, \texttt{tsum-mono},
\texttt{tsum-<-ε₀}; \Agda{BrouwerOrdinals.RectAlgebra}{\_⊴\_}.}

\begin{remark}
\label{rem:grading-rule}
One rule governs every construction below: an entry of the output is
raised only by entries of the input of its own order or higher. The
rule is what the iterates of $x \mapsto \omega^{x}$ break, since each
iterate raises the exponent by the value below it, and it is what
keeps the whole construction below~$\eps$, by
\cref{prop:closure}.
\end{remark}

\subsection{Scalar maps and cascades}
\label{sec:cascades}

At the ground type the maps needed are of the form
$a \mapsto (a \bplus t) \btimes \omega^{s}$, with a fixed
exponent~$s$. The combinator~$\Sc$ asks for
more, since the datum of $f\,a\,(g\,a)$ carries an exponent that varies
with~$a$, so we let the exponent be a map of the same kind one order
up.
\AgdaRefs{\Agda{BrouwerOrdinals.Move}{Move}; \texttt{img},
\texttt{orbit}, \texttt{move-orbit}, \texttt{padMove}.}

\begin{definition}[Cascades]
\label{def:cascade}
A \emph{cascade} of width~$k$ is a tuple of $k+1$ codes, one per
order of the datum it reads and a last one standing alone, and its
action on a datum of width~$k$ is defined by
\[
  \act{(c,E)}{(a,t)} = \big((a \bplus \tsum t) \bplus c\big)
    \btimes \omega^{\act{E}{t}} , \qquad
  \act{c}{\star} = c ,
\]
so that it reads the whole datum in its base and only the tail in its
exponent, the exponent being again a cascade. A cascade is
\emph{small} when each of its codes is below~$\eps$.
\AgdaRefs{\Agda{BrouwerOrdinals.Cascade}{Sc}; \texttt{valSc},
\texttt{ScAll<}.}
\end{definition}

\begin{remark}
\label{rem:last-grade}
The recursion ends at a free constant rather than at a translation,
and the choice is forced. A cascade whose exponent at the last grade
were read off a grade below it has no grade below to read, so that
exponent would be~$\bZ$; since $\omega^{\bZ} = \bS\,\bZ$ and
$x \btimes \bS\,\bZ = x$, the last grade would act by a translation,
and a translation cannot absorb a doubling: with $x$ the entry and $t$
the translation, $x \bplus x \le x \bplus t$ fails at
$x = t \bplus 1$, whatever $t$ is, since $x$ is chosen after it. With
a constant that the cascade carries itself, a doubling is paid at
every grade by one successor on the grade above.
\AgdaRefs{\Agda{BrouwerOrdinals.Cascade}{valSc-join};
\texttt{valSc-suc}.}
\end{remark}

\begin{lemma}
\label{lem:cascade-basic}
A small cascade takes a small datum to a code below~$\eps$, it is
monotone in the datum, its value at a datum is above the first entry
of the datum and above its own first code, and the sum of the values
of two cascades at one datum is below the value of a single cascade
there.
\AgdaRefs{\Agda{BrouwerOrdinals.Cascade}{valSc-<-ε₀};
\texttt{valSc-mono},
\texttt{valSc-join}, \texttt{valSc-increasing}, \texttt{valSc-const}.}
\end{lemma}

\subsection{Rectangular maps}
\label{sec:rect}

\begin{definition}[Graded maps]
\label{def:rect}
A \emph{graded map} from width~$p$ to width~$q$ is given, for each
output order~$j$, by a cascade reading the input orders~$j$ and above.
Its action is
\[
  \act{(E,D)}{(a,t)} = \big(\act{E}{(a,t)} , \act{D}{t}\big) ,
\]
the first component being a cascade on the whole input and the rest a
graded map applied to the tail. The datum of such a map is held as the
triangular array it is, one cascade per output order, each shorter
than the last, and a bare constant once the input orders are
exhausted.
\AgdaRefs{\Agda{BrouwerOrdinals.Rect}{RD}; \texttt{rect},
\texttt{RDAll<},
\texttt{rect-All<}, \texttt{rect-increasing}.}
\end{definition}

For the datum of a map to be a datum in its turn it has to be a tuple
as well. Writing $\W(p,q) = \max(p+1,q)$, a tuple $d$ of width
$\W(p,q)$ is read as a graded map $\rd d$ from width~$p$ to
width~$q$, the cascade for output order~$j$ being the suffix of the
tuple from~$j$. The sharing of entries is what makes the datum of a
map fit into a width one more than that of the input and no more than
that of the output.
\AgdaRefs{\Agda{BrouwerOrdinals.Rect}{WD}; \texttt{rdOf},
\texttt{rdOf-All<}, \texttt{rdOf-mono}.}

\begin{lemma}[The calculus]
\label{lem:calculus}
Graded maps support application, abstraction, projection and
iteration. Write
$y \sqcup x$ for the combination of a datum $y$ of width~$p$ with a
datum $x$ of width~$n$ into one of width $\max(p,n)$ containing both.
\begin{enumerate}
\item\label{item:apply}
  Application. For maps $F$ from width~$m$ to width $\W(p,q)$ and $G$
  from width~$m$ to width~$p$ there is a map $F\cdot G$ from width~$m$
  to width~$q$ with
  \[
    \act{\rd{\act F x}}{\act G x} \tle \act{(F\cdot G)}{x}
  \]
  for every datum~$x$ of width~$m$: a map applied to the result of
  another map of the same argument is a map of that argument.
\item\label{item:curry}
  Abstraction. For a map $D$ from width $\max(p,n)$ to width~$q$ there
  is a map $\Lambda D$ from width~$n$ to width $\W(p,q)$ with
  \[
    \act{D}{y \sqcup x} \tle \act{\rd{\act{(\Lambda D)}{x}}}{y}
  \]
  for all $y$ of width~$p$ and $x$ of width~$n$: a map of a combined
  argument is read back as a map of one part whose output acts on the
  other.
\item\label{item:var}
  Projection. If $t \tsub u$ then $t \tle \act{\mathrm{zero}}{u}$,
  where $\mathrm{zero}$ is the zero map, so a variable is bounded by
  its context.
\item\label{item:orbit}
  Iteration. Let $e$ be a datum of width $\W(p,p)$, let $x$ have
  width~$p$ and let $u$ be a small datum of width~$n$ with $e \tsub u$
  and $x \tsub u$. There is a map~$R$ from width~$n$ to width~$p$,
  depending only on those two widths, with
  \[
    \rd e^{\,j}(x) \tle \act{R}{u} \qquad \text{for every } j ,
  \]
  so the iterates of a square map applied to a datum inside a context
  are bounded by an action of that context, uniformly in the
  iteration count.
\end{enumerate}
In each case the resulting map is small when the given ones are.
\AgdaRefs{\Agda{BrouwerOrdinals.RectApply}{rect-ssub}; \texttt{sh},
\texttt{sh-All<}; \Agda{BrouwerOrdinals.RectCurry}{rect-curry},
\texttt{curryD}, \texttt{cmb};
\Agda{BrouwerOrdinals.RectAlgebra}{rect-⊴}, \texttt{zeroRD},
\texttt{flatRD}; \Agda{BrouwerOrdinals.RectOrbitAct}{rorbit-act},
\texttt{rorbitD}.}
\end{lemma}

Part~(\labelcref{item:orbit}) is what makes the bound on the recursor
a datum rather than an existential tuple, since the bound on the
iterates is produced once and for all from the context. The remaining
operations on these maps, the zero map, the join, the widening, the
flattening of a map into a tuple, the sum of the entries of a map and
the translation that the recursor applies after a supremum, are
routine and are collected in the formalization.
\AgdaRefs{\Agda{BrouwerOrdinals.RectAlgebra}{joinRD}; \texttt{padRD},
\texttt{addRDc}, \texttt{rect-addνc};
\Agda{BrouwerOrdinals.RectSum}{sumRD}, \texttt{tsum-rect};
\Agda{BrouwerOrdinals.RectAdd}{addν}, \texttt{rect-addν};
\Agda{BrouwerOrdinals.Absorb}{ι⊗ω\^{}}.}

\section{The bound on the majorant}
\label{sec:bound}

We define the predicate that a datum dominates a majorant and prove
the closure properties of it that the constants share, give the clause
for each constant, and conclude.

\subsection{The predicate}
\label{sec:predicate}

\begin{definition}[The order of a type and the domination predicate]
\label{def:dom}
The order of a type and the width of its data are
\[
  \ord(\gr) = 0 , \qquad
  \ord(\sigma\arr\tau) = \max(\ord(\sigma)+1, \ord(\tau)) ,
\]
\[
  \dwd(\sigma) = \ord(\sigma)+1 .
\]
A datum $d$ of width $\dwd(\sigma)$ \emph{dominates} a majorant $x$ of
type~$\sigma$, written $x \dom_{\sigma} d$, by recursion on the type:
at the ground type $a \dom_{\gr} d$ means $a \le \tsum d$, and
$F \dom_{\sigma\arr\tau} d$ means that
$x \dom_{\sigma} e$ implies $F\,x \dom_{\tau} \act{d}{e}$ for every
small~$e$, where $d$ is read as a graded map from width $\dwd(\sigma)$
to width $\dwd(\tau)$.
\AgdaRefs{\Agda{DialogueTrees.Predicate}{ord}; \texttt{wd},
\texttt{Dom}.}
\end{definition}

\begin{lemma}[The widths agree]
\label{lem:widths}
For all types $\sigma$ and~$\tau$,
$\dwd(\sigma\arr\tau) = \W(\dwd(\sigma), \dwd(\tau))$.
\AgdaRefs{The two sides are definitionally equal, which is what lets
the arrow clause of \Agda{DialogueTrees.Predicate}{Dom} typecheck
with no coercion, so there is no separate lemma to cite.}
\end{lemma}

\begin{proof}
Both sides are $\max(\ord(\sigma)+1, \ord(\tau))+1$, by the clause
$\max(m+1,n+1) = \max(m,n)+1$.
\end{proof}

The grading of the types and the grading of the data therefore agree
exactly, and the clause for an arrow reads the datum as a map and
applies it with no padding.

\begin{lemma}[Closure properties]
\label{lem:dom-closure}
Let $\sigma$ be a type.
\begin{enumerate}
\item\label{item:dom-mono}
  The predicate is monotone: if $x \dom_{\sigma} d$ and
  $d \tle d'$ then $x \dom_{\sigma} d'$.
\item\label{item:dom-sup}
  It is closed under pointwise suprema: if $h\,n \dom_{\sigma} d$ for
  every~$n$, then $\sup_{\sigma} h \dom_{\sigma} d$.
\item\label{item:dom-add}
  It is closed under the translation: if $x \dom_{\sigma} d$ and
  $\nu \le \nu'$, then
  $x \bplus_{\sigma} \nu \dom_{\sigma} d'$, where $d'$ adds $\nu'$ at
  order~$0$ and takes a successor at every order above.
\item\label{item:dom-eps}
  At the ground type a small datum gives a bound below~$\eps$: if
  $a \dom_{\gr} d$ and $d$ is small then $a < \eps$.
\end{enumerate}
\AgdaRefs{\Agda{DialogueTrees.Predicate}{Dom-mono};
\texttt{Dom-sup},
\texttt{Dom-⊕}, \texttt{Dom-<-ε₀}.}
\end{lemma}

\begin{proof}
Each is an induction on the type whose arrow case passes under the
argument. For (\labelcref{item:dom-add}), a translation alone would
not do: absorbing a code into a translation under a limit multiplier
asks for $b \btimes m \bplus d \le (b \bplus d) \btimes m$, which is
false, so the code is added at order~$0$ and a successor is taken at
every order above.
\end{proof}

\subsection{The clauses}
\label{sec:clauses}

\begin{lemma}[The easy clauses]
\label{lem:easy-clauses}
Any datum dominates the majorants of $\Zeroc$, $\Succc$ and~$\Kc$, and
the datum $(\bS\,\bZ, s)$ dominates the majorant of the oracle for
every~$s$. Application is the clause for an arrow read forwards.
\AgdaRefs{\Agda{DialogueTrees.Predicate}{Zero-clause};
\texttt{Succ-clause}, \texttt{Ω-clause}, \texttt{K-clause},
\texttt{app-clause}.}
\end{lemma}

\begin{proof}
The clause for $\Zeroc$ is the least element property of~$\bZ$, and
the one for $\Succc$ holds because a datum is below its image under a
graded map. The oracle adds one query and therefore
asks that the slot of order~$0$ be positive, which
$\bS\,\bZ$ supplies. For~$\Kc$, whose majorant discards its second
argument, the datum of the first argument has to survive two
applications; the output of the first application still dominates the
datum it was given, and the output of the second still dominates the
datum it was read from, and composing the two under monotonicity is
the whole proof.
\AgdaRefs{\Agda{BrouwerOrdinals.Rect}{rect-increasing};
\texttt{rect-arg-≥}, \texttt{rect-rdOf-≥}.}
\end{proof}

\begin{lemma}[The clause for $\Sc$]
\label{lem:s-clause}
At every triple of types there is a small datum that dominates the
majorant of~$\Sc$, depending only on the three widths.
\AgdaRefs{\Agda{DialogueTrees.SClause}{dS}; \texttt{dS-All<},
\texttt{S-clause}.}
\end{lemma}

\begin{proof}
The datum is built rather than designed, from
\cref{lem:calculus}. The proof is the one a reader of
$\lambda f\,g\,a.\,f\,a\,(g\,a)$ would write: the three data arrive one
at a time and are combined, in the order in which they are bound, into
one datum; the body is two applications over that one datum, which is
\cref{lem:calculus}(\labelcref{item:apply}) twice; and the three
binders are~\cref{lem:calculus}(\labelcref{item:curry}) twice and the
flattening of a map into a tuple, which is abstraction with an empty
context.
\end{proof}

\begin{lemma}[The clause for the recursor]
\label{lem:iter-clause}
At every type there is a small datum that dominates the majorant of
the recursor.
\AgdaRefs{\Agda{DialogueTrees.IterClause}{dI}; \texttt{dI-All<},
\texttt{Iter-clause}.}
\end{lemma}

\begin{proof}
As in~\cref{lem:s-clause}. The bound on the iterates is supplied by
iteration, \cref{lem:calculus}(\labelcref{item:orbit}), and the
translation by the count is absorbed by the operation of
\cref{lem:dom-closure}(\labelcref{item:dom-add}), the code added
belonging to the context as well.
\end{proof}

\subsection{The theorem}
\label{sec:theorem}

\begin{theorem}
\label{thm:fundamental}
Every closed term $t$ of $\TmO\sigma$ has a small datum of width
$\dwd(\sigma)$ dominating its majorant.
\AgdaRefs{\Agda{DialogueTrees.Fundamental}{fundamental}.}
\end{theorem}

\begin{proof}
By induction on the term, the constants by~\cref{lem:easy-clauses},
\cref{lem:s-clause} and~\cref{lem:iter-clause}, and application by the
clause for an arrow, whose datum is the image of the datum of the
argument under the datum of the function read as a graded map.
\end{proof}

\begin{corollary}
\label{cor:mu-below}
Every closed term $t$ of $\TmO\gr$ has $\Mj(t) < \eps$.
\AgdaRefs{\Agda{DialogueTrees.Fundamental}{μᵀ-<-ε₀}.}
\end{corollary}

\begin{proof}
At the ground type the datum is a single code, so
\cref{lem:dom-closure}(\labelcref{item:dom-eps}) applies.
\end{proof}

\begin{theorem}[The upper bound]
\label{thm:main}
For every closed term $t : (\gr\arr\gr)\arr\gr$ of system T,
\[
  \height(\dtree(t)) < \eps .
\]
\AgdaRefs{\Agda{DialogueTrees.Fundamental}{dialogue-tree-height-<-ε₀};
\Agda{DialogueTrees.Main}{restated}.}
\end{theorem}

\begin{proof}
\Cref{cor:mu-below} discharges the hypothesis of
\cref{thm:reduction}.
\end{proof}

\section{Sharpness}
\label{sec:sharp}

We exhibit terms of small height, read a dialogue tree off every
Brouwer code, and construct for each~$n$ a term attaining the
approximant $\tower n$, from which $\eps$ is the least bound.

\subsection{Terms of small height}
\label{sec:small}

\begin{examples}
\label{ex:small}
The three terms
$\lambda f.\,0$, $\lambda f.\,f\,0$ and
$\lambda f.\,\iterop\,f\,(f\,0)\,(f\,0)$ have dialogue trees
$\eta\,0$, $\beta\,\eta\,0$ and
$\beta(j \mapsto \iterop\,\mathrm{generic}\,(\beta\,\eta\,0)\,j)\,0$,
each by computation. The first has height~$\bZ$, and the second has
height~$\bS\,\bZ$ in the sense that both bounds hold, the code
$\bL(n \mapsto \bS\,\bZ)$ being the one the definition returns.
\AgdaRefs{\Agda{DialogueTrees.Examples}{leaf-term};
\texttt{query-term}, \texttt{iterated-term}, \texttt{leaf-tree},
\texttt{query-tree}, \texttt{iterated-tree}, \texttt{leaf-height},
\texttt{query-height-≤}, \texttt{query-height-≥}.}
\end{examples}

\begin{proposition}
\label{prop:omega-attained}
The third term of~\cref{ex:small} has height~$\omega$, which is
$\tower 0$.
\AgdaRefs{\Agda{DialogueTrees.Examples}{iterated-height-≥};
\texttt{iterated-height-≤}.}
\end{proposition}

\begin{proof}
By the computation of the tree in~\cref{ex:small}, the height is the
supremum over~$j$ of one more than the height of the $j$-th iterate of
the generic oracle at $\beta\,\eta\,0$. If the numeral $\num k$ is
below the height of a tree, then $\num{k+1}$ is below the height of
its image under the generic oracle, by induction on the tree, the
numeral being carried through a single branch at each step as
in~\cref{rem:one-branch}. So $\num j$ is below the height of the
$j$-th iterate, which gives the lower bound. In the other direction
the $j$-th iterate has height at most $\num 1 \bplus \num j$, which
is~$\num{j+1}$, so every branch is below~$\omega$.
\AgdaRefs{\Agda{DialogueTrees.Examples}{ι-generic};
\texttt{ι-iter},
\texttt{iter-generic-≤}.}
\end{proof}

\subsection{Trees of arbitrary height}
\label{sec:trees}

\begin{proposition}
\label{prop:tree-of}
Every Brouwer code is below the height of a dialogue tree read off the
code. In particular $\eps$ is below the height of such a tree, and
that tree is the dialogue tree of no term of system T.
\AgdaRefs{\Agda{DialogueTrees.Examples}{tree-of};
\texttt{code-≤-height}, \texttt{approximant}, \texttt{epsilon₀},
\texttt{no-term-of-height-ε₀}.}
\end{proposition}

\begin{proof}
A successor becomes a query whose subtrees are all the tree of the
predecessor, and a limit becomes a query whose $n$-th subtree is the
tree of the $n$-th member, and $a \le \height(\mathrm{tree}(a))$
follows by induction on~$a$. Were a term~$t$ to have the tree of~$\eps$,
then $\eps \le \height(\dtree(t))$, so that
$\bS\,\eps \le \bS\,\height(\dtree(t)) \le \eps$ by monotonicity of
the successor and \cref{thm:main}, and $\bS\,\eps \le \eps$ is refuted
by induction on the code.
\AgdaRefs{\Agda{DialogueTrees.Examples}{S-irrefl}.}
\end{proof}

The bound of~\cref{thm:main} is thus a statement about the terms, the
trees themselves reaching every height.

\subsection{Terms attaining each approximant}
\label{sec:church}

The terms of this section are Church encodings of ordinals, following
Coquand, Hancock and Setzer~\cite{CoquandHancockSetzer1997}; the same
device is applied to dialogue trees, which are Church encoded within
system T, in~\cite{EscardoPaivaRahliTosun2025}.

\begin{definition}[Church ordinals in system T]
\label{def:church}
A Church ordinal over a type~$\sigma$ folds itself over an
interpretation of zero, successor and limit, so it inhabits
\[
  \mathcal{O}\sigma = \sigma \arr (\sigma\arr\sigma)
    \arr ((\gr\arr\sigma)\arr\sigma) \arr \sigma .
\]
Zero, successor, limit and addition keep the parameter fixed, while
multiplication and exponentiation consume an encoding one level up,
exponentiation having the type
$\mathcal{O}(\mathcal{O}\sigma) \arr \mathcal{O}(\mathcal{O}\sigma)
 \arr \mathcal{O}\sigma$. The approximant $\tower n$ is therefore
written at its own type, with $n$ nestings of~$\mathcal{O}$ below the
level at which it is used.
\AgdaRefs{\Agda{DialogueTrees.Sharpness}{𝓞}; \texttt{zerᵀ},
\texttt{sucᵀ}, \texttt{limᵀ}, \texttt{addᵀ}, \texttt{mulᵀ},
\texttt{expᵀ}, \texttt{towerᵀ}; \Agda{Ordinals.ChurchEncoding}{O}.}
\end{definition}

\begin{remark}
\label{rem:no-eps}
Every approximant of~$\eps$ is written in system T in this way, while
$\eps$ itself calls for a uniform type and stays out of reach, in
accordance with~\cref{thm:main}. Dependent types and a universe give
the uniform typing and hence~$\eps$, which is how
\cite{CoquandHancockSetzer1997} obtains it.
\end{remark}

\begin{definition}[The terms]
\label{def:tower-term}
The term $t_{n} : (\gr\arr\gr)\arr\gr$ runs the encoding of $\tower n$
at the ground type against the oracle~$\alpha$ through the algebra
that interprets zero as~$0$, the successor of~$x$ as
$\iterop\,(\lambda y.\,y)\,x\,(\alpha\,0)$, and the limit of a
sequence~$\varphi$ as $\varphi\,(\alpha\,0)$.
\AgdaRefs{\Agda{DialogueTrees.Sharpness}{tower-term}; \texttt{abs},
\texttt{tm}.}
\end{definition}

In the standard interpretation the successor of the algebra is the
identity, since iterating the identity returns~$x$ whatever the count.
What the dialogue interpretation makes of it is another matter,
because the recursor reads its count through a Kleisli extension, so
the tree of the count is placed above the tree of~$x$. The successor
therefore places one query above the tree built so far: the oracle is
asked at~$0$ and every branch of the answer returns~$x$, so the
height of the result exceeds the height of~$x$ by one. An
interpretation of the successor that asks the oracle below, as the
term of~\cref{prop:omega-attained} does, raises every leaf by one and
leaves the height of an infinite tree where it was, which is why the
query goes on top. The limit hands its query to the sequence, and the
sequence branches over the answer, with the tree of the $j$-th member
on the $j$-th branch.

\begin{theorem}[The lower bound]
\label{thm:lower}
For every~$n$,
\[
  \tower n \le \height(\dtree(t_{n})) .
\]
\AgdaRefs{\Agda{DialogueTrees.Sharpness}{tower-attained};
\Agda{DialogueTrees.Main}{attained}.}
\end{theorem}

\begin{proof}
By parametricity. The relation $\mathrm{Rel}(\sigma, R)$ states that a
Church ordinal sends every algebra realising the three code
constructors to a realisation of the fold of a Brouwer code, where
realisation at the parameter type is an arbitrary relation~$R$ and the
limit of the algebra is fed Kleisli extensions alone, which is what
the terms produce and what turns one query into a branch per index.
Zero, successor, limit, addition, multiplication and exponentiation
preserve the relation, moving the code along the matching operation on
codes, and each proof is a direct unfolding, the dialogue semantics of
a compiled closed term computing. Instantiating the relation at the
ground type with ``the code is below the height of the tree'', the
three conditions of the algebra are checked directly. The code side
lands on a fold of~$\tower n$, which is $\tower n$ itself, the folds
used for the product and the power grouping their summands in the
other order and agreeing with the operations
of~\cref{def:arithmetic} by associativity of addition.
\AgdaRefs{\Agda{DialogueTrees.Sharpness}{Rel}; \texttt{fold},
\texttt{mulᶜ}, \texttt{expᶜ}, \texttt{towerᶜ-tower},
\texttt{expᶜ-ω\^{}}, \texttt{towerᵀ-Rel}.}
\end{proof}

\begin{corollary}
\label{cor:least}
The heights of the dialogue trees of closed terms of type
$(\gr\arr\gr)\arr\gr$ are cofinal in~$\eps$, and $\eps$ is the least
code bounding all of them.
\AgdaRefs{\Agda{DialogueTrees.Sharpness}{heights-cofinal};
\texttt{ε₀-least}, \texttt{tower-term-below-ε₀};
\Agda{DialogueTrees.Main}{cofinal}, \texttt{least}.}
\end{corollary}

\begin{proof}
A code below~$\eps$ is below some $\tower n$, which
\cref{thm:lower} attains. If $b$ bounds the height of every tree then
$\tower n \le b$ for every~$n$, so $\eps \le b$ by the limiting
property of the order.
\end{proof}

\section{Discussion}
\label{sec:discussion}

We place the bound and the method in the literature, and compare the
proof with that of Xiao and Sun.

\subsection{The ordinal and the method}
\label{sec:method}

The bound $\eps$ is the proof-theoretic ordinal of arithmetic, by
Gentzen~\cite{Gentzen1936}; a modern account is
Pohlers~\cite{Pohlers2009}. The hereditary assignment of ordinals to
terms, type by type, is the method of Howard~\cite{Howard1970}, who
bounds the values a term computes and the length of its reductions.
The dialogue-tree height is a different measure, and
\cref{thm:main} is not a corollary of that assignment, but the shape
of the induction is the same.

The present proof has two features independent of the result. The
first is that the ordinals are syntactic. Brouwer codes
with the order of~\cref{def:codes} make the successor monotone by
construction, which is what allows every lemma to be proved without
excluded middle; a treatment through the ordinals of the univalent
library needs it, and the development records what fails. The second
is the grading. The predicate of~\cref{def:dom} is indexed by the
order of a type, and~\cref{lem:widths} says that this grading is the
same as the grading of the data, so that the clause for an arrow
applies a graded map with no adjustment. Once that is in place the two
hard clauses are written in the calculus of~\cref{sec:graded} rather
than by hand.

\subsection{The proof of Xiao and Sun}
\label{sec:xiaosun}

The bound of \cref{thm:main} is due to Xiao and
Sun~\cite{XiaoSun2026}, and the proof given here is an alternative
one. Their argument runs through the syntax. A recursor is translated
into a closed infinitary template, the redexes are removed by a
sequence of passes indexed by ordinary type level, the translation and
every pass preserve the dialogue denotation, and an auxiliary rank
with an additive substitution bound is tracked across the passes. The
outcome is sharper than \cref{thm:main} in the respect that matters
for a given term: they compute from the term a natural
number~$K(t) \ge 2$ with
\[
  \height(\dtree(t)) < \tower{K(t)} < \eps ,
\]
so each term is bounded by a named level of the tower rather than
by~$\eps$. The proof here names no such level.

The two proofs differ in their ordinals. Theirs are the classical
ones, and the ordinal part of their formalization assumes excluded
middle for the propositions of the second universe, together with
univalence, propositional truncation and set replacement, while their
syntactic part assumes nothing. We have been unable to locate their
Agda formalization for comparison with ours. The present development
replaces the ordinals by Brouwer codes with the syntactic order
throughout, so that no part of it is classical and function
extensionality is the only assumption. Since the syntactic order is
sound for the standard interpretation of the codes, \cref{thm:main}
gives the classical statement as well.

They write that they do not claim their bound sharp, and leave open
whether the heights are cofinal in~$\eps$. \Cref{cor:least} answers
that question: they are, so $\eps$ is the least bound, and no tower
level bounds every term.

\begin{question}
\label{q:exact}
The exact height of a term as a function of its syntax is determined
by neither proof. Nor is the least~$n$ with
$\height(\dtree(t)) < \tower n$ for a given term~$t$, for which the
number $K(t)$ of~\cite{XiaoSun2026} is an upper bound.
\end{question}

\appendix
% amsart keeps the section counter for appendices, so cleveref has to
% be told that a section label after \appendix names an appendix.
\crefalias{section}{appendix}
\crefalias{subsection}{appendix}

\section{The formalization}
\label{sec:appendix}

The formalization was developed first, by LLM's, and the paper was
then unformalized from it, again by LLM's, each statement of the paper
being written from the Agda statement it points at. We describe the
files of the development, say where the main theorems are and what
has to be read by hand, describe the pointers that tie the paper to
the code, and record how the paper was verified against it.

\subsection{The files}
\label{sec:files}

The development is 28 Agda modules under the prefix
\AgdaModule{DialogueTreesHeight}, some three thousand lines of code, in
two directories, available at
\url{https://github.com/martinescardo/DialogueTreesHeight}. Every
module is checked with \texttt{--safe} and \texttt{--without-K}, with
no postulates and no holes, and takes
function extensionality as a parameter named \texttt{fe}. The library
it builds on is \texttt{TypeTopology}: the Brouwer codes come from
\AgdaModule{Ordinals.BrouwerCodes}, and the dialogue trees, the
combinatory system T and its dialogue interpretation from
\AgdaModule{EffectfulForcing.MFPSAndVariations}.

The directory \AgdaModule{BrouwerOrdinals}, 17 modules, holds the
ordinal arithmetic and the graded maps, and knows nothing of system T
or of dialogue trees. The modules \AgdaModule{Order},
\AgdaModule{Multiplication}, \AgdaModule{Epsilon0} and
\AgdaModule{MultClosure} are \cref{sec:codes}: the order, the
arithmetic, the tower and the closure properties of~$\eps$. The other
thirteen are \cref{sec:graded}: \AgdaModule{MoveTuple} for the data,
\AgdaModule{Cascade} and \AgdaModule{CascadeLevel} for the cascades,
\AgdaModule{Rect} for the graded maps and the reading of a tuple as
one, \AgdaModule{RectApply}, \AgdaModule{RectCurry},
\AgdaModule{RectAlgebra} and \AgdaModule{RectOrbitAct} for the four
facts of \cref{lem:calculus}, and \AgdaModule{Move},
\AgdaModule{MultAffine}, \AgdaModule{Absorb}, \AgdaModule{RectAdd} and
\AgdaModule{RectSum} for the arithmetic those rest on.

The directory \AgdaModule{DialogueTrees}, 10 modules, is the rest
of the paper. The module \AgdaModule{Constructive} defines the height,
\cref{def:height}; \AgdaModule{Majorant} and \AgdaModule{MajorantAll}
are \cref{sec:majorant}; \AgdaModule{Predicate}, \AgdaModule{SClause},
\AgdaModule{IterClause} and \AgdaModule{Fundamental} are
\cref{sec:bound}; \AgdaModule{Examples} and \AgdaModule{Sharpness} are
\cref{sec:sharp}; and \AgdaModule{Main} restates the main theorems.
The module \AgdaModule{index} imports all of them in dependency order
and ends with an outline of the proof and the references.

\subsection{The main theorems}
\label{sec:main-theorems}

The results of the paper are the following four, in two modules of
the directory \AgdaModule{DialogueTrees}, and all four are restated
independently in its module \AgdaModule{Main}.
\begin{center}
\begin{tabular}{lll}
  \Cref{thm:main}, upper bound &
    \AgdaModule{Fundamental} & \texttt{dialogue-tree-height-<-ε₀} \\
  \quad restated &
    \AgdaModule{Main} & \texttt{restated} \\
  \Cref{thm:lower}, lower bound &
    \AgdaModule{Sharpness} & \texttt{tower-attained} \\
  \quad restated &
    \AgdaModule{Main} & \texttt{attained} \\
  \Cref{cor:least}, cofinality &
    \AgdaModule{Sharpness} & \texttt{heights-cofinal} \\
  \quad restated &
    \AgdaModule{Main} & \texttt{cofinal} \\
  \Cref{cor:least}, least bound &
    \AgdaModule{Sharpness} & \texttt{ε₀-least} \\
  \quad restated &
    \AgdaModule{Main} & \texttt{least}
\end{tabular}
\end{center}
Agda checks that each is inhabited. What it does not check is that
the statements mean what the paper says they mean, so these are what
has to be read by hand: the four types above, the type
\texttt{T ((ι ⇒ ι) ⇒ ι)} they quantify over, and the definitions they
are built from, which are the order and the arithmetic of
\cref{def:codes,def:arithmetic} in \AgdaModule{Order},
\AgdaModule{Multiplication} and \AgdaModule{Epsilon0}, the height of
\cref{def:height} in \AgdaModule{Constructive}, and
\texttt{dialogue-tree} in \texttt{TypeTopology}. Together they come to
some twenty lines. The module \AgdaModule{Main} guards the statements
against drift. It restates the upper bound with $<$ unfolded and
checks that the same proof term inhabits the restatement, checks that
the order is not degenerate at the bottom and that the type quantified
over is inhabited, and names the types of \texttt{dialogue-tree} and
\texttt{height}. It does the same for the three sharpness results,
with $<$ unfolded to a successor on the left and $\eps$ to the limit of
the tower, and checks by \texttt{refl} that the tower is the sequence
$\omega$, $\omega^{\omega}$, $\omega^{\omega^{\omega}}, \ldots$ and
that $\eps$ is its limit.

Two claims in the paper that something is false, in
\cref{rem:last-grade} and in the proof of
\cref{lem:dom-closure}(\labelcref{item:dom-add}), are motivation and
are not formalized, and the account of the proof of Xiao and Sun
in~\cref{sec:xiaosun} comes from reading their paper.

\subsection{The pointers}
\label{sec:deformalizing}

Every numbered item carries a pointer, set in blue, to the Agda that
proves it. A pointer names a module and a definition in it, then
further names of the same module, and a name must be defined in the
module it is cited under rather than re-exported by it. A script
resolves every pointer against the sources and reports, for this
paper, 52 notes carrying 199 names, all resolving. Four numbered items
carry no pointer, three remarks and a question, which make no formal
claim, and \cref{lem:widths} carries a note in place of a lemma, its
two sides being definitionally equal. The paper was written statement
first, each Agda statement being read before the paper's was written,
and the proofs describe the route the Agda takes.

\subsection{The A/B/C verification}
\label{sec:abc}

The paper was then verified against the formalization in three
passes, the independence of the first two being the point of the
arrangement.

\begin{enumerate}
\item[A.]
  Read the paper alone, with the formalization closed, as a reader
  with no access to it: whether the mathematics is correct, whether
  each statement is self-contained, whether each proof proves what
  its statement says in steps that can be followed, and whether the
  prose says what it means. The findings go to a numbered list.
\item[B.]
  Audit the same text against the formalization, without opening the
  list from pass~A: each statement against the Agda it points at,
  each proof against the route the Agda takes, and each pointer
  against what exists and in which module.
\item[C.]
  Read the two lists together, merge the findings that coincide, drop
  a finding of~A that B refutes, and rank what survives with the
  errors first, the faithfulness gaps next and the language last.
  Each finding is then settled before the next is raised.
\end{enumerate}

On the first run pass~A produced sixteen findings and pass~B eleven, one
pair merged, and twenty-five survived into pass~C, of which five were
errors. They are worth naming as a sample of what each pass catches.
Pass~A found an operation used in the statement of a lemma and
nowhere defined, a constant used in four places and nowhere defined,
and a definition that reads as trivially the identity unless one is
told why it is not. Pass~B found a proof attributing two of its cases
to a lemma that neither case uses, a second proof describing a route
the Agda does not take, and a theorem stated in a shape weaker than
the one the pointer proves.

The procedure was run more than once, and with more than one model,
the later runs finding no error of mathematics and a small number of
faithfulness gaps and points of language, most of them in text the
first run had added. A limitation of every run should be recorded.
The value of pass~A lies in the reader's ignorance of the
formalization, and in each run the model carrying out that pass had
the formalization at hand, so that it could read the paper as it was
meant rather than as it stands. An independent pass~A is the single
change that would most strengthen the procedure.

\begin{thebibliography}{9}

\bibitem{CoquandHancockSetzer1997}
Thierry Coquand, Peter Hancock and Anton Setzer.
\newblock Ordinals in type theory.
\newblock Invited talk, Computer Science Logic (CSL), 1997.
\newblock \url{https://www.cse.chalmers.se/~coquand/ordinal.pdf}

\bibitem{Escardo2013}
Mart\'in H\"otzel Escard\'o.
\newblock Continuity of G\"odel's system T definable functionals via
  effectful forcing.
\newblock {\em Electronic Notes in Theoretical Computer Science},
  298:119--141, 2013. MFPS XXIX.
\newblock \url{https://doi.org/10.1016/j.entcs.2013.09.010}

\bibitem{EscardoPaivaRahliTosun2025}
Mart\'in H\"otzel Escard\'o, Bruno da Rocha Paiva, Vincent Rahli and
  Ayberk Tosun.
\newblock Internal effectful forcing in system T.
\newblock In {\em 10th International Conference on Formal Structures
  for Computation and Deduction (FSCD 2025)}, volume 337 of {\em
  LIPIcs}, pages 19:1--19:17, 2025.
\newblock \url{https://doi.org/10.4230/LIPIcs.FSCD.2025.19}

\bibitem{Gentzen1936}
Gerhard Gentzen.
\newblock Die Widerspruchsfreiheit der reinen Zahlentheorie.
\newblock {\em Mathematische Annalen}, 112:493--565, 1936.
\newblock English translation in M.~E. Szabo, editor, {\em The
  Collected Papers of Gerhard Gentzen}, North-Holland, 1969.

\bibitem{Howard1970}
William A. Howard.
\newblock Assignment of ordinals to terms for primitive recursive
  functionals of finite type.
\newblock In A.~Kino, J.~Myhill and R.~E. Vesley, editors, {\em
  Intuitionism and Proof Theory (Buffalo 1968)}, pages 443--458.
  North-Holland, 1970.

\bibitem{Pohlers2009}
Wolfram Pohlers.
\newblock {\em Proof Theory: The First Step into Impredicativity}.
\newblock Universitext. Springer, 2009.

\bibitem{XiaoSun2026}
MingKun Xiao and YiXuan Sun.
\newblock An explicit ordinal bound for system T dialogue trees.
\newblock Preprint, arXiv:2609.20369, 2026.
\newblock \url{https://arxiv.org/abs/2609.20369}

\end{thebibliography}

\end{document}
